% locking/locking-existence.tex
% mainfile: ../perfbook.tex
% SPDX-License-Identifier: CC-BY-SA-3.0

\section{基于锁的存在性保证}
\label{sec:locking:Lock-Based Existence Guarantees}
%
\epigraph{存在先于本质并支配本质。}{Jean-Paul Sartre}

并行编程中的一个关键挑战是提供
\emph{\IXpl{existence guarantee}}~\cite{Gamsa99}，
以便对给定对象的访问尝试可以依赖于该对象在整个访问尝试期间一直存在。

\begin{listing}
\begin{fcvlabel}[ln:locking:Per-Element Locking Without Existence Guarantees]
\begin{VerbatimL}[commandchars=\\\@\$]
int delete(int key)
{
	int b;
	struct element *p;

	b = hashfunction(key);
	p = hashtable[b];
	if (p == NULL || p->key != key)		\lnlbl@chk_first$
		return 0;
	spin_lock(&p->lock);			\lnlbl@acq$
	hashtable[b] = NULL;			\lnlbl@NULL$
	spin_unlock(&p->lock);
	kfree(p);
	return 1;				\lnlbl@return1$
}
\end{VerbatimL}
\end{fcvlabel}
\caption{无存在性保证的逐元素加锁（有缺陷！）}
\label{lst:locking:Per-Element Locking Without Existence Guarantees (Buggy!)}
\end{listing}

在某些情况下，存在性保证是隐式的：

\begin{enumerate}
\item	基础模块中的全局变量和静态局部变量将在应用程序运行期间一直存在。
\item	已加载模块中的全局变量和静态局部变量将在该模块保持加载状态期间一直存在。
\item	只要模块的某个函数有活跃实例，该模块就会保持加载状态。
\item	给定函数实例的栈上变量将一直存在，直到该实例返回。
\item	如果你正在某个函数中执行，或者被该函数直接或间接调用，
	那么该函数就有一个活跃实例。
\end{enumerate}

这些隐式存在性保证很直观，但涉及隐式存在性保证的
bug 确实可能发生。

\QuickQuiz{
	依赖隐式存在性保证怎么会导致 bug？
}\QuickQuizAnswer{
	以下是一些因不当使用隐式存在性保证而导致的 bug：
	\begin{enumerate}
	\item	程序将一个全局变量的地址写入文件，
		然后该程序的后续实例读取该地址并尝试解引用它。
		由于地址空间随机化，这可能会失败，
		更不用说程序的重新编译了。
	\item	一个模块可以将其某个变量的地址记录在
		位于其他模块中的指针中，然后在该模块被
		卸载后尝试解引用该指针。
	\item	一个函数可以将其某个栈上变量的地址记录到
		一个全局指针中，其他函数可能在该函数
		返回后尝试解引用该指针。
	\end{enumerate}
	我相信你能想出更多的可能性。
}\QuickQuizEnd

但更有趣——也更棘手——的保证涉及堆内存：
动态分配的数据结构将一直存在，直到它被释放。
要解决的问题是如何将该结构的释放与对同一结构的并发访问同步。
一种方法是使用\emph{显式保证}，例如加锁。
如果只有在持有某个锁时才能释放给定结构，那么持有该锁就保证了该结构的存在。

但这个保证依赖于锁本身的存在。
保证锁存在的一种直接方法是将锁放在全局变量中，但全局加锁的缺点是限制了可扩展性。
一种随数据结构大小增长而提高可扩展性的方法是在结构的每个元素中放置一个锁。
不幸的是，将用于保护数据元素的锁放在数据元素本身中
容易产生微妙的\IXpl{race condition}，如
\cref{lst:locking:Per-Element Locking Without Existence Guarantees (Buggy!)}
所示。

\QuickQuiz{
	\begin{fcvref}[ln:locking:Per-Element Locking Without Existence Guarantees]
	如果我们需要删除的元素不是
	\cref{lst:locking:Per-Element Locking Without Existence Guarantees (Buggy!)}
	\clnref{chk_first} 上列表的第一个元素怎么办？
	\end{fcvref}
}\QuickQuizAnswer{
	这是一个没有链接的非常简单的哈希表，所以给定桶中唯一的
	元素就是第一个元素。
	读者可以尝试将此示例改编为具有完整链接的哈希表。
}\QuickQuizEnd


\begin{fcvref}[ln:locking:Per-Element Locking Without Existence Guarantees]
要看到其中一个 race condition，请考虑以下事件序列：
\begin{enumerate}
\item	线程~0 调用 \co{delete(0)}，执行到代码清单的 \clnref{acq}，获取了锁。
\item	线程~1 并发调用 \co{delete(0)}，也执行到 \clnref{acq}，
	但因线程~0 持有锁而自旋等待。
\item	线程~0 执行 \clnrefrange{NULL}{return1}，将元素从
	哈希表中移除，释放锁，然后释放该元素。
\item	线程~0 继续执行，并分配内存，恰好得到了
	它刚刚释放的那块内存。
\item	线程~0 然后将这块内存初始化为
	其他类型的结构。
\item	线程~1 的 \co{spin_lock()} 操作失败，因为
	它认为是 \co{p->lock} 的东西已经不再
	是一个 spinlock 了。
\end{enumerate}
因为没有存在性保证，当线程试图在 \clnref{acq} 上获取
该元素的锁时，数据元素的身份可能已经改变了！
\end{fcvref}

\begin{listing}
\begin{fcvlabel}[ln:locking:Per-Element Locking With Lock-Based Existence Guarantees]
\begin{VerbatimL}[commandchars=\\\@\$]
int delete(int key)
{
	int b;
	struct element *p;
	spinlock_t *sp;

	b = hashfunction(key);
	sp = &locktable[b];
	spin_lock(sp);				\lnlbl@acq$
	p = hashtable[b];			\lnlbl@getp$
	if (p == NULL || p->key != key) {
		spin_unlock(sp);
		return 0;
	}
	hashtable[b] = NULL;
	spin_unlock(sp);
	kfree(p);
	return 1;
}
\end{VerbatimL}
\end{fcvlabel}
\caption{具有基于锁的存在性保证的逐元素加锁}
\label{lst:locking:Per-Element Locking With Lock-Based Existence Guarantees}
\end{listing}

\begin{fcvref}[ln:locking:Per-Element Locking With Lock-Based Existence Guarantees]
修复此示例的一种方法是使用一组哈希化的全局锁，使每个哈希桶都有自己的锁，如
\cref{lst:locking:Per-Element Locking With Lock-Based Existence Guarantees}
所示。
这种方法允许在获取数据元素的指针之前（在 \clnref{getp}）先获取适当的锁（在 \clnref{acq}）。
虽然这种方法对于包含在单个可分区数据结构（如代码清单所示的哈希表）中的元素效果很好，
但如果给定数据元素可以是多个哈希表的成员，
或者对于更复杂的数据结构如树或图，这种方法可能会有问题。
这些问题不仅可以解决，而且其解决方案还构成了基于锁的软件事务内存
实现的基础~\cite{Shavit95,DaveDice2006DISC}。
然而，
\cref{chp:Deferred Processing}
描述了提供存在性保证的更简单——也更快——的方法。
\end{fcvref}

% @@@ Optimized sharded locking from P1726R5, and pointer-zap implications.
